
\section{Methodische Vorgehensweise}
In diesem Kapitel soll gezeigt werden, wie die Übersetzung der Notebooks erfolgt sind. Von dem zuerst grundlegenden Verständnis über die einzelnen Funktionen hin zu der unterstützenden Übersetzung mithilfe von \ac{KI} bis hin zum vollständig funktionierenden Notebook

\subsection{Github Projects}
Hier möchte ich unser Projektstatus zeigen, eingehen, wie wir uns diese Tickets zugeteilt haben und wie wir mit der großen Menge an Notebooks fertig geworden sind.


\subsection{Auswahl der Lehrmaterialien}
Welche Kriterien es gab...


\subsection{Der Übersetzungs-Workflow}
Wann manuell, wann automatisiert


\subsection{Einsatz von KI-gestützten Tools}
Verwendete Modelle (GPT-4, Claude, Copilot etc.)\\
Prompt-Engineering-Strategien für Code-Übersetzung. (hier würde ich gerne unseren Promt erwähnen und erläutern, vielleicht machen wir hierfür nochmal ein Unterkapitel)\\
Grenzen der KI (Halluzinationen bei seltenen TS-Libraries).